\scanpage{145}
caractérisera $\mathcal M$ à équivalence près). De plus, on aime donner des conditions moyennant lesquelles le foncteur $\underline h$ est \uline{pleinement fidèle}.

On va supposer que la structure \deleted{\unsure{pré}}additive sur $\mathscr V$ peut s'enrichir en une structure $k$-linéaire, où $\deleted{\unsure{k}}Q$ est un \uline{corps} commutatif,
% slanted
\margin{i.e. on a un hom d'anneaux $Q\to\operatorname{End}\,\mathrm{id}_{\mathscr V}$}
[NB. si $Q=\mathbb Q$, cet enrichissement est déterminé de façon unique]. On suppose que \deleted{les \unsure{Hom} $\operatorname{Fl}(X,Y)$ sont des \unsure{esp vectoriels|espaces vectoriels} de dim finie sur $k$, et} que pour tt $X$, $\operatorname{Fl}(X,X)$ est une \uline{algèbre} \inserted{\uline{de degré fini}} \uline{semi-simple} sur $\deleted{\unsure{k}}Q$. Nous allons alors \deleted{\unsure{de|en}} construire une catégorie $\mathcal M^{\deleted{\unsure{\circ|0|\#}}}$ de la façon suivante. Un objet de $\mathcal M^{\deleted{\circ}}$ est un couple $(X,\pi)$, où $X\in\operatorname{Ob}\mathscr V$, $\pi\in\operatorname{Fl}(X,X)$, $\pi^2=\pi$ ($\pi$ est un \uline{projecteur} de $X$). Un homomorphisme $(X,\pi)\to(X',\pi')$ est un hom $f:X\to X'$ tel que $\pi'f=f\pi$, \uline{modulo} les $f$ tels que $f\pi=\pi'f=0$.
\diagram{\begin{tikzcd}[ampersand replacement=\&,row sep=small,column sep=small]
X\arrow[r,dash,"\pi"]\arrow[d,"f"'] \& X\arrow[d,"f"]\\
X'\arrow[r,dash,"\pi'"]\arrow[d,"g"'] \& X'\arrow[d,"g"]\\
X''\arrow[r,dash,"\pi''"] \& X''
\end{tikzcd}}
On compose de façon évidente. D'où une catégorie\inserted{ $\mathcal M^{\deleted{\circ}}$}, manifestement additive, et un foncteur additif
% archive sheet number: 144
